\chapter{Implementation}
\label{chap:implementation}

This chapter describes the data pipeline, the modifications made to the base
implementation, and the defects found while making them.

\section{Environment}
\label{sec:environment}

The base implementation ships no dependency versions. Packages were pinned to
the releases current at the base paper's submission: torch 2.3.1, torchvision
0.18.1, numpy 1.26.4, pandas 2.2.2. This is not conservatism for its own sake.
The optimiser module uses a positional-alpha overload of \texttt{Tensor.add\_}
that later torch releases reject, so the Per-FedAvg baseline fails outright on a
current install.

All experiments run on CPU. Threading was measured rather than assumed: the
default of one thread per core is slower than a single thread for these model
sizes, by twenty-two to thirty-nine per cent, because the matrix products are
too small for the synchronisation cost to pay. Runs are therefore pinned to one
thread and parallelised across configurations instead.

\section{Data pipeline}
\label{sec:pipeline}

\subsection{CICIDS2017}

Two distributions of CICIDS2017 exist. The machine learning variant carries
seventy-eight features and a label; the traffic labelling variant adds flow
identifier, addresses, ports and timestamp. Only the second supports a temporal
split, and only the second supports the host-based partition, so it is the one
used. This distinction matters beyond this project: a temporal split performed
on the first variant can only be positional, since it holds no time column.

Cleaning proceeds in four steps. One file contains 288,602 rows in which every
field is empty; removing them leaves a total of 2,830,743, matching the
canonical figure. Seven of the eight files record timestamps in twelve-hour
form with no meridiem, so the Wednesday capture appears to run from 08:42 to
02:43; since the capture window is 08:00 to 17:00, hours one to seven are
afternoon and are advanced by twelve. Rows carrying infinities or missing values
in the flow rate columns are dropped along with exact duplicates, 309,745 rows
in total. Identifier columns are removed before the feature matrix is built.

Classes below a configurable threshold are dropped. At 2,000 rows this removes
Bot, the three web attack variants, Infiltration and Heartbleed, leaving nine
classes. The threshold exists because a class with eleven rows in the entire
capture cannot survive a seventy-five/twenty-five split, and its F1 contribution
is noise that enters the macro average with equal weight.

Feature values are clipped to percentile bounds fitted on the training split
and then standardised, also on training statistics only. Clipping was added
after measuring the range: fifty-eight of seventy-nine features have a
maximum-to-median ratio above $10^4$ and forty-three above $10^6$.

\subsection{NSL-KDD and TON-IoT}

NSL-KDD ships a canonical train and test split, so no split has to be designed
and the questions raised in Section~\ref{sec:eval-protocol} do not arise. Its
three categorical columns are one-hot encoded, giving 122 features over five
classes: normal plus the four documented attack domains.

TON-IoT contributes its network subset, 211,043 rows over ten classes. It is
included because it is near-balanced where CICIDS2017 is not: normal traffic is
23.7 per cent against 81.7, so the accuracy floor is 0.2369 rather than 0.8171
and macro F1 has room to move. Addresses are dropped as features, since the
attacker hosts are a fixed range and would leak the label, but the destination
address is retained to define the host-based partition.

\section{Metric instrumentation}
\label{sec:metrics-impl}

The base implementation records accuracy and loss only. No confusion matrix is
constructed anywhere in it, and the published paper reports no F1, recall or
precision.

A metrics module was added computing macro-averaged F1, precision, recall and
false positive rate from a pooled confusion matrix. Two decisions in it matter.

The average is taken over the complete label set, with classes absent from a
test slice contributing zero. Averaging only over classes that happen to appear
inflates personalised methods, because each device holds a subset of classes
and would otherwise be graded only on those.

The matrices are summed across clients before the F1 is taken. Averaging
per-client F1 scores instead produces a different and more flattering number,
for the same reason.

Confusion accumulation is vectorised through \texttt{bincount}. The obvious
Python loop over prediction pairs dominates the round loop at these test set
sizes.

\section{Modifications to the base implementation}
\label{sec:modifications}

Four changes were made, each defaulting to the shipped behaviour.

\subsection{Argument binding}

\texttt{UserPerMFL} passed five positional arguments into the six parameter
slots that follow \texttt{beta} in the base class, shifting every subsequent
binding by one. The \texttt{--lamda} argument was bound to \texttt{self.eta} and
\texttt{--local\_iters} to \texttt{self.lamda}. The optimiser therefore received
a penalty weight of twenty, the local iteration count, rather than the value
supplied on the command line. The per-step pull toward the team model was 0.2
where the paper's own setting gives 0.005, a factor of forty.

The correction is keyword arguments. Since \texttt{self.eta},
\texttt{self.gamma} and \texttt{self.local\_epochs} are assigned but never read
in this class, the penalty weight is the only behavioural change.

\subsection{Seeding}

\texttt{torch.manual\_seed(0)} was fixed in the entry point, the loaders fix
their own seeds, and client selection seeds on the round index. Under sequential
team assignment nothing varied between repeats, so the repeat count produced
identical runs and no standard deviation could be reported. This is consistent
with the standard deviations of exactly zero reported throughout the base
paper's main results table. A seed argument now drives all three.

\subsection{Decoupled team weight}

Equations~\eqref{eq:device} and~\eqref{eq:team} both contain $\lambda$.
Section~\ref{sec:lambda-analysis} shows the two roles want opposed values. A
second argument sets the team-update weight independently, defaulting to the
device value so that unmodified behaviour is preserved.

\subsection{Team formation from client updates}

Teams are assigned in the data loader, before any client has trained, so no
update-based grouping can be computed there. A clustering module was added that
runs in the server's round loop instead. The signal is the proximal residual
$\theta_{i,j} - w_i$, which the device update already computes. It measures how
far a device would move given the pull its team exerts, which neither of the
comparable methods can use, since neither has a proximal term. Formation follows
SCMoE-PFL's procedure: L2 normalisation, PCA, cosine similarity, centres seeded
from the least similar pairs. Assignment is hard and size-constrained rather
than soft and overlapping, because the server sizes its team arrays statically
and divides by per-team sample counts.

\section{Defects in the base implementation}
\label{sec:defects}

Sixteen defects were catalogued. Those that affect results are listed here; the
remainder appear in Appendix~\missing{ref}.

\begin{description}
  \item[Argument binding.] As above. Silently overrides a documented
    hyperparameter.
  \item[Non-reproducible partitioning.] Random team assignment reseeds on the
    wall clock and prints the seed, so teams differ on every run at a fixed
    torch seed. The clock seed then propagates into per-round team selection.
  \item[Unseeded synthetic split.] The synthetic loader seeds numpy but not
    Python's random module before the per-client shuffle, so the split differs
    between processes. Verified by hashing across three runs.
  \item[Broken partitioning mode.] One of the three team assignment modes
    assigns the return value of an in-place shuffle, which is \texttt{None},
    then indexes it. That mode has never executed.
  \item[Dataset mislabelling.] The EMNIST path loads the sixty-two class
    by-class split, which the paper distinguishes as a separate dataset. The
    paper states the digits split was used.
  \item[Unreachable branch.] A reshape guard tests a dataset name that differs
    in case from the value the parser accepts, so that dataset skips the
    reshape and its convolutional path cannot run.
  \item[Silent no-op.] The experiment counter is both the loop start and
    compared against the repeat count, so a start index at or above the repeat
    count performs no work and exits zero.
  \item[Result overwriting.] Output filenames encode the hyperparameters but
    not the round count, client count or local iteration count, so runs
    differing only in those overwrite each other.
  \item[Unreachable algorithm.] The parser accepts an algorithm name that the
    dispatch does not match, so selecting it raises a name error. It is one of
    the base paper's headline comparison methods.
  \item[Broken aggregation helper.] The routine that averages results across
    repeats is commented out at its call site and passes two arguments to a
    three-argument definition.
  \item[Missing regularisation.] Appendix D.3 describes multinomial logistic
    regression with L2 regularisation. No weight decay or penalty term appears
    anywhere in the repository. This matters for the strongly convex
    convergence result, whose step size bound depends on a strong convexity
    modulus that is zero without it.
  \item[Model depth.] Sections 4 and D.3 describe networks with two hidden
    layers. The implementation has one. Both were measured; the paper's
    description performs worse on this data by 0.011 to 0.021 macro F1.
\end{description}
